\section{Introduction}

\section{GG-KM}

The GG-KM (an acronym for Generalized Gamma Distribution and Kernel
Machines) is a novel semiparametric approach to long-term survival
modeling. The main objective is to address the limitations of linear
predictors by employing kernel functions to model
$\log\left(\theta(\mathbf{x})\right)$ within a Reproducing Kernel Hilbert
Space (RKHS), where $\theta(\mathbf{x})$ is a parameter governing the
distribution of the latent number of competing causes in the promotion
cure model. This formulation allows for a flexible, potentially
nonlinear relationship between the covariates and the latent competing-risk
structure.

To provide greater flexibility in modeling the survival distribution of
susceptible individuals, the Generalized Gamma distribution is adopted for
the latent event times. This distribution includes several commonly used
parametric survival distributions as special cases and, consequently,
allows a broad range of hazard shapes to be represented.

The model is estimated through an Expectation--Maximization (EM)
algorithm and is formulated for different probability distributions of the
latent number of competing causes $N$, namely the Bernoulli, Binomial,
Poisson, and Negative Binomial distributions. This formulation provides a
unified framework for investigating different assumptions regarding the
latent competing-risk structure while preserving the kernel-based
semiparametric formulation.

\subsection{Nonparametric extension via kernel methods}

Although the classical promotion-time cure model allows covariates to be
incorporated through parametric functions of $\theta(\mathbf{x})$, such an
approach may be restrictive when the relationship between the covariates
and the latent competing-risk structure is complex or highly nonlinear.

In this context, the proposition is to model the mean intensity of the latent
competing causes, $\theta(\mathbf{x})$, through a nonparametric function
\begin{equation}
f(\mathbf{x}) = \log\left(\theta(\mathbf{x})\right),
\end{equation}
where $f:\mathcal{X}\rightarrow\mathbb{R}$. The logarithmic
transformation ensures that $\theta(\mathbf{x})>0$ for every
$\mathbf{x}$. The estimation of $f$ is formulated through the minimization
of a penalized negative log-likelihood:
\begin{equation}
\label{eq:ggkm_penalized_likelihood}
\min_{f \in \mathcal{H}_K}
\left\{
-\sum_{i=1}^{n} l_i(f)
+
\frac{\lambda}{2}\|f\|_{\mathcal{H}_K}^{2}
\right\},
\end{equation}
where $l_i(f)$ denotes the log-likelihood contribution of the $i$th
individual, $\mathcal{H}_K$ is a Reproducing Kernel Hilbert Space (RKHS)
associated with a positive-definite kernel $K(\cdot,\cdot)$, and
$\lambda\geq 0$ is a regularization parameter controlling the complexity
of the estimated function.

By the Representer Theorem, the solution to the optimization problem in
Equation~\eqref{eq:ggkm_penalized_likelihood} admits a finite-dimensional
representation of the form
\begin{equation}
\label{eq:ggkm_representer}
f^*(\mathbf{x})
=
\sum_{i=1}^{n}\alpha_i K(\mathbf{x},\mathbf{x}_i),
\end{equation}
where $\alpha_1,\ldots,\alpha_n\in\mathbb{R}$ are coefficients to be
estimated. Consequently, the mean intensity of the latent competing
causes for an individual with covariate vector $\mathbf{x}_i$ is given by
\begin{equation}
\label{eq:ggkm_theta}
\theta(\mathbf{x}_i)
=
\exp\left\{
\sum_{j=1}^{n}
\alpha_j K(\mathbf{x}_i,\mathbf{x}_j)
\right\}.
\end{equation}

\subsection{Estimação via algoritmo EM penalizado}

A representação obtida pelo Teorema do Representante permite reescrever o modelo de tempo de promoção em
termos de um número finito de parâmetros. Seja
\begin{equation}
f(\mathbf{x}_i) = \sum_{j=1}^{n} \alpha_j K(\mathbf{x}_i,\mathbf{x}_j),
\end{equation}

Sob a formulação do modelo de tempo de promoção, assume-se que o número de riscos latentes associados ao
indivíduo $i$ segue uma distribuição de Poisson com média $\theta(\mathbf{x}_i)$, isto é,
\begin{equation}
N_i
\sim
\text{Poisson}
\left(
\theta(\mathbf{x}_i)
\right).
\end{equation}

Para realizar a estimação dos parâmetros, foi considerado o algoritmo EM. Em particular, considera-se uma extensão
particular, o EM penalizado \cite{mclachlan2008algorithm}.

% Entretanto, os riscos competitivos $N_i$ associados ao evento de interesse não são observados. Assim, a
% estimação dos parâmetros pode ser conduzida por meio do algoritmo EM. Neste caso, considera-se o algoritmo
% EM penalizado \cite{mclachlan2008algorithm}.

Dessa forma, considere os dados completos $D_{\text{comp}} = \left\{t_i, \delta_i, N_i, \mathbf{x}_i\right\}^n_{i=1}$,
em que $t_i$ representa o tempo observado, $\delta_i$ o indicador de falha e $N_i$ o número de riscos latentes
associados ao $i\text{-ésimo}$ indíviduo. A log-verossimilhança completa, menos uma constante aditiva
$C = -\sum_{i=1}^n \log\left(N_i!\right)$, pode ser escrita como
\begin{equation}
\begin{aligned}
l_c
=
&
\sum_{i=1}^{n}
\left[
-\exp\!\bigl(f(\mathbf{x}_i)\bigr)
+
N_i f(\mathbf{x}_i)
\right]
\\
&
+
\sum_{i=1}^{n}
\left[
\delta_i
\log f_{GG}(t_i;a,d,p)
\right]
\\
&
+
\sum_{i=1}^{n}
\left[
(N_i-\delta_i)
\log S_{GG}(t_i;a,d,p)
\right].
\end{aligned}
\end{equation}

Portanto, tem-se a log-verossimilhança completa penalizada:

\begin{equation}
\ell_{c, P} = l_c - \frac{\lambda}{2} \boldsymbol{\alpha}^{\top} K \boldsymbol{\alpha},
\end{equation}
em que
\begin{equation}
\boldsymbol{\alpha} = (\alpha_1,\ldots,\alpha_n)^{\top}
\end{equation}
e $K$ denota a matriz de kernel com elementos $K_{ij}=K(\mathbf{x}_i,\mathbf{x}_j)$.

O algoritmo EM consiste em alternar entre o cálculo das esperanças condicionais das variáveis latentes e a maximização da esperança da log-verossimilhança completa penalizada.

No passo E, as quantidades não observadas são substituídas por suas respectivas esperanças condicionais, calculadas sob os valores correntes dos parâmetros. Dessa forma, define-se a função auxiliar

\begin{equation}
Q\!\left(
\Theta \mid \Theta^{(k)}
\right)
=
\mathbb{E}
\left[
\ell_{c,P}
\mid
\mathcal{D}_{\mathrm{obs}},
\Theta^{(k)}
\right],
\end{equation}

em que
\[
\Theta=(\boldsymbol{\alpha},a,d,p)
\]
representa o vetor de parâmetros do modelo.

Utilizando propriedades do modelo de tempo de promoção, obtém-se que a esperança condicional do número de riscos latentes possui forma fechada, dada por

\begin{equation}
\widehat{N}_i^{(k+1)}
=
\mathbb{E}
\!\left[
N_i
\mid
\mathcal{D}_{\mathrm{obs}},
\Theta^{(k)}
\right]
=
\delta_i
+
\theta_i^{(k)}
S_{GG}
\!\left(
t_i;
a^{(k)},
d^{(k)},
p^{(k)}
\right),
\end{equation}

em que

\begin{equation}
\theta_i^{(k)}
=
\exp
\!\left(
f^{(k)}(\mathbf{x}_i)
\right)
=
\exp
\!\left(
\sum_{j=1}^{n}
\alpha_j^{(k)}
K(\mathbf{x}_i,\mathbf{x}_j)
\right).
\end{equation}

Substituindo \(N_i\) por sua esperança condicional na log-verossimilhança completa penalizada e eliminando
termos constantes em relação aos parâmetros, obtém-se a função auxiliar penalizada
\begin{equation}
Q_P
=
Q_{\alpha}
+
Q_{GG},
\end{equation}
em que
\begin{equation}
Q_{\alpha}
=
\sum_{i=1}^{n}
\left[
-\exp\!\bigl(f(\mathbf{x}_i)\bigr)
+
\widehat N_i^{(k+1)}
f(\mathbf{x}_i)
\right]
-
\frac{\lambda}{2}
\boldsymbol{\alpha}^{\top}
K
\boldsymbol{\alpha},
\end{equation}
e
\begin{equation}
Q_{GG}
=
\sum_{i=1}^{n}
\left[
\delta_i
\log f_{GG}(t_i;a,d,p)
+
\bigl(\widehat N_i^{(k+1)}-\delta_i\bigr)
\log S_{GG}(t_i;a,d,p)
\right].
\end{equation}

Observa-se que a função auxiliar penalizada é separável em dois blocos de parâmetros: o primeiro associado aos coeficientes do kernel, \(\boldsymbol{\alpha}\), e o segundo aos parâmetros da distribuição Gama Generalizada, \((a,d,p)\). Essa decomposição permite que o passo M seja realizado por meio de atualizações alternadas dos dois blocos.

Inicialmente, atualizam-se os coeficientes do kernel resolvendo

\begin{equation}
\boldsymbol{\alpha}^{(k+1)}
=
\arg\max_{\boldsymbol{\alpha}}
Q_{\alpha}.
\end{equation}

Em seguida, atualizam-se os parâmetros da distribuição Gama Generalizada por meio de
\begin{equation}
(a^{(k+1)},d^{(k+1)},p^{(k+1)})
=
\arg\max_{a,d,p}
Q_{GG}.
\end{equation}
% Cabe destacar que o termo de penalização afeta apenas a atualização dos coeficientes do kernel, não
% alterando a forma da atualização dos parâmetros da distribuição Gama Generalizada.
Ambos os problemas de maximização são resolvidos numericamente por meio do algoritmo L-BFGS-B, utilizando as derivadas analíticas
das respectivas funções objetivo.

Após cada iteração, calcula-se a log-verossimilhança penalizada observada
\begin{equation}
\ell_P
=
\sum_{i=1}^{n}
\left[
\delta_i
\left(
f(\mathbf{x}_i)
+
\log f_{GG}(t_i;a,d,p)
\right)
-
\exp\!\bigl(f(\mathbf{x}_i)\bigr)
F_{GG}(t_i;a,d,p)
\right]
-
\frac{\lambda}{2}
\boldsymbol{\alpha}^{\top}
K
\boldsymbol{\alpha},
\end{equation}
a qual é utilizada para monitorar a convergência do algoritmo. O procedimento iterativo é interrompido quando
a variação absoluta da log-verossimilhança penalizada observada entre duas iterações consecutivas se torna
inferior a uma tolerância previamente especificada, isto é,
\begin{equation}
\left|
\ell_P^{(k+1)}
-
\ell_P^{(k)}
\right|
<
\varepsilon,
\end{equation}
em que \(\varepsilon>0\) representa uma constante de tolerância.

% O algoritmo EM consiste em alternar entre o cálculo da esperança condicional dos riscos latentes e a
% maximização da esperança da log-verossimilhança completa penalizada.

% No passo E, calcula-se
% \begin{equation}
% Q\!\left(
% \Theta \mid \Theta^{(k)}
% \right)
% =
% \mathbb{E}
% \left[
% \ell_{c,P}
% \mid
% \mathcal{D}_{\mathrm{obs}},
% \Theta^{(k)}
% \right].
% \end{equation}
% em que $\Theta = \left(\boldsymbol{\alpha},a,d,p\right)$.

% Pode-se provar que:
% % Utilizando propriedades do modelo de tempo de promoção, obtém-se que
% \begin{equation}
% \widehat{N}_i^{(k+1)}
% =
% \mathbb{E}
% \!\left[
% N_i
% \mid
% \mathcal{D}_{obs},
% \Theta^{(k)}
% \right]
% =
% \delta_i
% +
% \theta_i^{(k)}
% S_{GG}
% \!\left(
% t_i;
% a^{(k)},
% d^{(k)},
% p^{(k)}
% \right),
% \end{equation}
% em que
% \begin{equation}
% \theta_i^{(k)}
% =
% \exp
% \!\left(
% f^{(k)}(\mathbf{x}_i)
% \right)
% =
% \exp
% \!\left(
% \sum_{j=1}^{n}
% \alpha_j^{(k)}
% K(\mathbf{x}_i,\mathbf{x}_j)
% \right).
% \end{equation}

% Substituindo as quantidades não observadas por suas respectivas
% esperanças condicionais, a função auxiliar penalizada assume a forma
% \begin{equation}
% Q_P
% =
% Q_{\alpha}
% +
% Q_{GG},
% \end{equation}
% com
% \begin{equation}
% Q_{\alpha}
% =
% \sum_{i=1}^{n}
% \left[
% -\exp\!\bigl(f(\mathbf{x}_i)\bigr)
% +
% \widehat N_i^{(k+1)}
% f(\mathbf{x}_i)
% \right]
% -
% \frac{\lambda}{2}
% \boldsymbol{\alpha}^{\top}
% K
% \boldsymbol{\alpha},
% \end{equation}
% e
% \begin{equation}
% Q_{GG}
% =
% \sum_{i=1}^{n}
% \left[
% \delta_i
% \log f_{GG}(t_i;a,d,p)
% +
% \bigl(\widehat N_i^{(k+1)}-\delta_i\bigr)
% \log S_{GG}(t_i;a,d,p)
% \right].
% \end{equation}

% Observa-se que a função auxiliar penalizada é separável em dois blocos de parâmetros,
% \(\boldsymbol{\alpha}\) e \((a,d,p)\). Consequentemente, o passo M pode ser realizado por
% otimização alternada.

% Inicialmente, atualizam-se os coeficientes do kernel resolvendo
% \begin{equation}
% \boldsymbol{\alpha}^{(k+1)}
% =
% \arg\max_{\boldsymbol{\alpha}}
% Q_{\alpha}.
% \end{equation}

% Em seguida, atualizam-se os parâmetros da distribuição Gama Generalizada por meio de
% \begin{equation}
% (a^{(k+1)},d^{(k+1)},p^{(k+1)})
% =
% \arg\max_{a,d,p}
% Q_{GG}.
% \end{equation}

% A otimização de ambos problemas de otimização é realizada numericamente através de L-BFGS-B, utilizando-se
% as derivadas analíticas das respectivas funções objetivos.

% O processo iterativo é repetido até a convergência da log-verossimilhança penalizada observada,
% \begin{equation}
% \ell_P
% =
% \sum_{i=1}^{n}
% \left[
% \delta_i
% \left(
% f(\mathbf{x}_i)
% +
% \log f_{GG}(t_i;a,d,p)
% \right)
% -
% \exp\!\bigl(f(\mathbf{x}_i)\bigr)
% F_{GG}(t_i;a,d,p)
% \right]
% -
% \frac{\lambda}{2}
% \boldsymbol{\alpha}^{\top}
% K
% \boldsymbol{\alpha},
% \end{equation}
% sendo adotado como critério de parada
% \begin{equation}
% \left|
% \ell_P^{(k+1)}
% -
% \ell_P^{(k)}
% \right|
% <
% \varepsilon,
% \end{equation}
% em que \(\varepsilon > 0\) representa uma tolerância previamente especificada.

\begin{algorithm}[H]
\caption{Algoritmo EM penalizado para estimar o modelo GG-KM - $N_i \sim \mathrm{Poisson}\!\bigl(\theta(\mathbf{x}_i)\bigr).$}
\begin{algorithmic}[1]

\State Inicialize
$
\alpha^{(0)},\,
a^{(0)},\,
d^{(0)},\,
p^{(0)}
$
e escolha uma tolerância $\varepsilon>0$.

\State Defina $k \gets 0$.

\Repeat

\State \textbf{E-step:} para $i=1,\ldots,n$, calcule
\[
\widehat N_i^{(k+1)}
=
\delta_i
+
\exp\!\left(
\sum_{j=1}^{n}
\alpha_j^{(k)}
K(x_i,x_j)
\right)
S_{GG}
\!\left(
t_i;
a^{(k)},
d^{(k)},
p^{(k)}
\right).
\]

\State \textbf{M-step:}
\[
\alpha^{(k+1)}
=
\arg\max_{\alpha}
\left\{
\sum_{i=1}^{n}
\left[
-
\exp\!\left(
\sum_{j=1}^{n}
\alpha_j
K(x_i,x_j)
\right)
+
\widehat N_i^{(k+1)}
\sum_{j=1}^{n}
\alpha_j
K(x_i,x_j)
\right]
-
\frac{\lambda}{2}
\alpha^\top K\alpha
\right\}.
\]

\[
(a^{(k+1)},d^{(k+1)},p^{(k+1)})
=
\arg\max_{a,d,p}
\sum_{i=1}^{n}
\left[
\delta_i
\log f_{GG}(t_i;a,d,p)
+
(\widehat N_i^{(k+1)}-\delta_i)
\log S_{GG}(t_i;a,d,p)
\right].
\]

\State Calcule a log-verossimilhança penalizada observada
\[
\ell_P^{(k+1)}
=
\ell_P
\!\left(
\alpha^{(k+1)},
a^{(k+1)},
d^{(k+1)},
p^{(k+1)}
\right).
\]

\State Atualize
\[
\Delta^{(k+1)}
=
\left|
\ell_P^{(k+1)}
-
\ell_P^{(k)}
\right|.
\]

\State $k \gets k+1$.

\Until{$\Delta^{(k+1)} < \varepsilon$}

\end{algorithmic}
\end{algorithm}

\subsection{Estimação do modelo GG-KM com riscos latentes Binomial Negativa}

Uma generalização do modelo apresentado anteriormente consiste em assumir que o número de riscos latentes
\(N_i\) segue uma distribuição Binomial Negativa em vez de uma distribuição de Poisson. Em vez de
especificar essa distribuição por meio de sua função de probabilidade, adota-se a parametrização definida
pela função geradora de probabilidades
\begin{equation}
A(u)
=
\left[
1+\phi\theta_i(1-u)
\right]^{-1/\phi},
\end{equation}
onde \(\theta_i>0\) representa o número médio de riscos latentes e \(\phi\) é um parâmetro de dispersão.
Sob essa parametrização,

\begin{equation}
\mathbb{E}(N_i)=\theta_i,
\qquad
\mathrm{Var}(N_i)
=
\theta_i+\phi\theta_i^2.
\end{equation}

Assim, valores positivos de \(\phi\) produzem sobredispersão em relação ao modelo de Poisson. Além disso,
essa parametrização engloba importantes casos particulares. Quando \(\phi \to 0\), recupera-se a
distribuição de Poisson e, consequentemente, o modelo GG-KM apresentado na seção anterior. Já o caso
\(\phi=-1\) conduz ao modelo de mistura de \citeonline{berkson1952survival}. Neste trabalho, entretanto,
considera-se apenas o cenário \(\phi>0\), correspondente à presença de sobredispersão.

A função de sobrevivência populacional passa a ser
\begin{equation}
S_{\text{pop}}(t)
=
\left[
1+\phi\theta_i
F_{GG}(t;a,d,p)
\right]^{-1/\phi},
\end{equation}
com
\begin{equation}
\theta_i
=
\exp
\!\left(
\sum_{j=1}^{n}
\alpha_j
K(\mathbf{x}_i,\mathbf{x}_j)
\right).
\end{equation}

Consequentemente, a log-verossimilhança penalizada observada é dada por

\begin{equation}
\ell_P
=
\sum_{i=1}^{n}
\left\{
\delta_i
\left[
f(\mathbf{x}_i)
+
\log
f_{GG}(t_i;a,d,p)
\right]
-
\left(
\frac{1}{\phi}
+
\delta_i
\right)
\log
\left[
1+
\phi
\theta_i
F_{GG}(t_i;a,d,p)
\right]
\right\}
-
\frac{\lambda}{2}
\boldsymbol{\alpha}^{\top}
K
\boldsymbol{\alpha}.
\end{equation}

Quando o parâmetro de dispersão \(\phi\) é previamente especificado, ele é mantido fixo ao
longo do processo iterativo e apenas os parâmetros \(\boldsymbol{\alpha}\), \(a\), \(d\) e \(p\) são estimados.
Por outro lado, quando \(\phi\) não é previamente informado, sua estimação é realizada conjuntamente com os
demais parâmetros do modelo.

Neste último caso, a estimação é realizada por meio de uma extensão do algoritmo EM, o
Expectation Conditional Maximization Either (ECME) \cite{mclachlan2008algorithm}. A escolha desse
algoritmo decorre da presença do parâmetro de dispersão \(\phi\), cuja
atualização não admite solução analítica em forma fechada. De fato,
\(\phi\) aparece simultaneamente nos termos \(\log(1+\phi\theta_iF_{GG}(t_i))\), \(1/\phi\) e nas
esperanças condicionais dos riscos latentes, o que torna a maximização da função auxiliar substancialmente
mais complexa do que a atualização dos demais parâmetros. Assim, os parâmetros \(\boldsymbol{\alpha}\)
e \((a,d,p)\) são atualizados por meio da maximização da função auxiliar, enquanto \(\phi\) é atualizado
diretamente pela maximização da log-verossimilhança observada.

% Neste último caso, a estimação é realizada por meio de uma extensão do algoritmo EM, o ECME
% (Expectation Conditional Maximization Either), \cite{mclachlan2008algorithm}. A escolha desse algoritmo decorre da presença do
% parâmetro de dispersão \(\phi\), cuja atualização a partir da função auxiliar não conduz a uma etapa de
% maximização simples. Assim, os parâmetros \(\boldsymbol{\alpha}\) e \((a,d,p)\) são atualizados por meio da
% maximização da função auxiliar, enquanto \(\phi\) é atualizado diretamente pela maximização da
% log-verossimilhança observada.

No passo E, obtém-se

\begin{equation}
\widehat N_i^{(k+1)}
=
\mathbb{E}
\!\left[
N_i
\mid
\mathcal D_{\mathrm{obs}},
\Theta^{(k)}
\right]
=
\delta_i
+
\frac{
\left(
1+\phi^{(k)}\delta_i
\right)
\theta_i^{(k)}
S_{GG}
\!\left(
t_i;
a^{(k)},
d^{(k)},
p^{(k)}
\right)
}{
1+
\phi^{(k)}
\theta_i^{(k)}
F_{GG}
\!\left(
t_i;
a^{(k)},
d^{(k)},
p^{(k)}
\right)
}.
\end{equation}

Substituindo as quantidades não observadas por suas respectivas esperanças condicionais, obtém-se uma
função auxiliar penalizada que pode ser decomposta em

\begin{equation}
Q_P
=
Q_{\alpha}
+
Q_{GG}.
\end{equation}

O termo associado aos parâmetros da distribuição Gama Generalizada, \(Q_{GG}\), permanece idêntico ao
obtido no caso Poisson. Já a componente dependente de \(\boldsymbol{\alpha}\) passa a ser

\begin{equation}
Q_{\alpha}
=
\sum_{i=1}^{n}
\left[
\widehat N_i^{(k+1)}
\log(\theta_i)
-
\left(
\widehat N_i^{(k+1)}
+
\frac{1}{\phi}
\right)
\log(1+\phi\theta_i)
\right]
-
\frac{\lambda}{2}
\boldsymbol{\alpha}^{\top}
K
\boldsymbol{\alpha}.
\end{equation}

O passo CM consiste então nas atualizações
\begin{equation}
\boldsymbol{\alpha}^{(k+1)}
=
\arg\max_{\boldsymbol{\alpha}}
Q_{\alpha},
\end{equation}
e
\begin{equation}
(a^{(k+1)},d^{(k+1)},p^{(k+1)})
=
\arg\max_{a,d,p}
Q_{GG}.
\end{equation}

Quando \(\phi\) não é previamente especificado, realiza-se adicionalmente a etapa ECME

\begin{equation}
\phi^{(k+1)}
=
\arg\max_{\phi>0}
\;
\ell_P
\!\left(
\boldsymbol{\alpha}^{(k+1)},
a^{(k+1)},
d^{(k+1)},
p^{(k+1)},
\phi
\right).
\end{equation}

A otimização dos problemas de maximização é realizada numericamente por meio do algoritmo L-BFGS-B,
utilizando-se as derivadas analíticas das respectivas funções objetivo. O processo iterativo é repetido
até a convergência da log-verossimilhança penalizada observada.

\begin{algorithm}[H]
\caption{Algoritmo ECME para estimar o modelo GG-KM com riscos latentes Binomial Negativa}
\begin{algorithmic}[1]

\State Inicialize
$
\alpha^{(0)},\,
a^{(0)},\,
d^{(0)},\,
p^{(0)},\,
\phi^{(0)}
$
e escolha uma tolerância $\varepsilon>0$.

\State Defina $k \gets 0$.

\Repeat

\State Para $i=1,\ldots,n$, calcule
\[
\theta_i^{(k)}
=
\exp\!\left(
\sum_{j=1}^{n}
\alpha_j^{(k)}
K(x_i,x_j)
\right),
\]
\[
F_{GG}^{(k)}(t_i)
=
F_{GG}
\!\left(
t_i;
a^{(k)},
d^{(k)},
p^{(k)}
\right),
\qquad
S_{GG}^{(k)}(t_i)
=
1-F_{GG}^{(k)}(t_i).
\]

\State \textbf{E-step:}
\[
\widehat N_i^{(k+1)}
=
\delta_i
+
\frac{
\bigl(1+\phi^{(k)}\delta_i\bigr)
\theta_i^{(k)}
S_{GG}^{(k)}(t_i)
}{
1+\phi^{(k)}
\theta_i^{(k)}
F_{GG}^{(k)}(t_i)
},
\qquad
i=1,\ldots,n.
\]

\State \textbf{M-step:}
\[
\alpha^{(k+1)}
=
\arg\max_{\alpha}
\left\{
\sum_{i=1}^{n}
\left[
\widehat N_i^{(k+1)}
\log(\theta_i)
-
\left(
\widehat N_i^{(k+1)}
+
\frac{1}{\phi^{(k)}}
\right)
\log\!\bigl(1+\phi^{(k)}\theta_i\bigr)
\right]
-
\frac{\lambda}{2}
\alpha^\top K\alpha
\right\},
\]

\[
(a^{(k+1)},d^{(k+1)},p^{(k+1)})
=
\arg\max_{a,d,p}
\sum_{i=1}^{n}
\left[
\delta_i
\log f_{GG}(t_i;a,d,p)
+
(\widehat N_i^{(k+1)}-\delta_i)
\log S_{GG}(t_i;a,d,p)
\right].
\]

\[
\phi^{(k+1)}
=
\arg\max_{\phi>0}
\,
\ell_O
\!\left(
\alpha^{(k+1)},
a^{(k+1)},
d^{(k+1)},
p^{(k+1)},
\phi
\right).
\]

\State Calcule a log-verossimilhança penalizada observada
\[
\ell_P^{(k+1)}
=
\ell_P
\!\left(
\alpha^{(k+1)},
a^{(k+1)},
d^{(k+1)},
p^{(k+1)},
\phi^{(k+1)}
\right).
\]

\State Atualize
\[
\Delta^{(k+1)}
=
\left|
\ell_P^{(k+1)}
-
\ell_P^{(k)}
\right|.
\]

\State $k \gets k+1$.

\Until{$\Delta^{(k+1)} < \varepsilon$}

\end{algorithmic}
\end{algorithm}

\begin{algorithm}[H]
\caption{Algoritmo EM penalizado para estimar o modelo GG-KM - $N_i \sim \mathrm{Binomial}(K,\theta_i^*)$.}
\begin{algorithmic}[1]

\State Inicialize
$
\alpha^{(0)},
a^{(0)},
d^{(0)},
p^{(0)}
$
e escolha uma tolerância $\varepsilon>0$.

\State Defina $k \gets 0$.

\Repeat

\State \textbf{E-step:} para $i=1,\ldots,n$, calcule

\[
f_i^{(k)}
=
\sum_{j=1}^{n}
\alpha_j^{(k)}
K(x_i,x_j),
\]

\[
\theta_i^{*(k)}
=
\frac{\exp(f_i^{(k)})}
     {1+\exp(f_i^{(k)})},
\]

e

\[
\widehat N_i^{(k+1)}
=
\delta_i
+
(K-\delta_i)
\frac{
\theta_i^{*(k)}
S_{GG}
\!\left(
t_i;
a^{(k)},
d^{(k)},
p^{(k)}
\right)
}
{
1-\theta_i^{*(k)}
+
\theta_i^{*(k)}
S_{GG}
\!\left(
t_i;
a^{(k)},
d^{(k)},
p^{(k)}
\right)
}.
\]

\State \textbf{M-step:}

Atualize $\alpha$ resolvendo

\[
\alpha^{(k+1)}
=
\arg\max_{\alpha}
\left\{
\sum_{i=1}^{n}
\left[
\widehat N_i^{(k+1)}
f_i
-
K
\log
\bigl(
1+\exp(f_i)
\bigr)
\right]
-
\frac{\lambda}{2}
\alpha^\top K\alpha
\right\},
\]

onde

\[
f_i
=
\sum_{j=1}^{n}
\alpha_j
K(x_i,x_j).
\]

\State Atualize os parâmetros da distribuição Generalizada Gama:

\[
(a^{(k+1)},d^{(k+1)},p^{(k+1)})
=
\arg\max_{a,d,p}
\sum_{i=1}^{n}
\left[
\delta_i
\log f_{GG}(t_i;a,d,p)
+
(\widehat N_i^{(k+1)}-\delta_i)
\log S_{GG}(t_i;a,d,p)
\right].
\]

\State Calcule a log-verossimilhança penalizada observada

\[
\ell_P^{(k+1)}
=
\ell_P
\!\left(
\alpha^{(k+1)},
a^{(k+1)},
d^{(k+1)},
p^{(k+1)}
\right).
\]

\State Atualize

\[
\Delta^{(k+1)}
=
\left|
\ell_P^{(k+1)}
-
\ell_P^{(k)}
\right|.
\]

\State $k \gets k+1$.

\Until{$\Delta^{(k+1)} < \varepsilon$}

\end{algorithmic}
\end{algorithm}

\begin{algorithm}[H]
\caption{EM with functional gradient boosting for the GG-KM model -
$N_i \sim \mathrm{Binomial}(K,\theta_i^*)$.}
\begin{algorithmic}[1]

\State Initialize
$
f^{(0)}(\cdot),
a^{(0)},
d^{(0)},
p^{(0)}
$
and choose $\varepsilon>0$, $\varepsilon_{\mathrm{tree}}>0$, and
the maximum number of trees $M$.

\State Set $k\gets0$.

\Repeat

\State \textbf{E-step:} compute, for $i=1,\ldots,n$,

\[
\theta_i^{*(k)}
=
\frac{\exp\{f^{(k)}(x_i)\}}
     {1+\exp\{f^{(k)}(x_i)\}},
\]

\[
\widehat N_i^{(k+1)}
=
\delta_i
+
(K-\delta_i)
\frac{
\theta_i^{*(k)}
S_{GG}(t_i;a^{(k)},d^{(k)},p^{(k)})
}{
1-\theta_i^{*(k)}
+
\theta_i^{*(k)}
S_{GG}(t_i;a^{(k)},d^{(k)},p^{(k)})
}.
\]

\State \textbf{M-step for $f(x)$:} set $f_0(x)\gets f^{(k)}(x)$ and,
for $m=1,\ldots,M$, compute

\[
g_{i,m}
=
\widehat N_i^{(k+1)}
-
K
\frac{\exp\{f_{m-1}(x_i)\}}
     {1+\exp\{f_{m-1}(x_i)\}},
\]

fit a regression tree $h_m(x)$ to
$\{(x_i,g_{i,m})\}_{i=1}^n$, and obtain

\[
\rho_m
=
\arg\max_{\rho\geq0}
\sum_{i=1}^n
\left[
\widehat N_i^{(k+1)}
\{f_{m-1}(x_i)+\rho h_m(x_i)\}
-
K\log
\left(
1+\exp\{f_{m-1}(x_i)+\rho h_m(x_i)\}
\right)
\right].
\]

\State Update

\[
f_m(x)=f_{m-1}(x)+\rho_m h_m(x),
\]

and stop the boosting iterations if

\[
\frac{1}{n}\sum_{i=1}^n
|f_m(x_i)-f_{m-1}(x_i)|
<\varepsilon_{\mathrm{tree}}.
\]

\State Set $f^{(k+1)}(x)\gets f_m(x)$.

\State \textbf{M-step for the Generalized Gamma parameters:}

\[
(a^{(k+1)},d^{(k+1)},p^{(k+1)})
=
\arg\max_{a,d,p}
\sum_{i=1}^n
\left[
\delta_i\log f_{GG}(t_i;a,d,p)
+
(\widehat N_i^{(k+1)}-\delta_i)
\log S_{GG}(t_i;a,d,p)
\right].
\]

\State Compute

\[
\Delta^{(k+1)}
=
\left|
\ell_P^{(k+1)}-\ell_P^{(k)}
\right|,
\]

where

\[
\ell_P^{(k+1)}
=
\ell_P
\left(
f^{(k+1)},a^{(k+1)},d^{(k+1)},p^{(k+1)}
\right).
\]

\State $k\gets k+1$.

\Until{$\Delta^{(k+1)}<\varepsilon$}

\end{algorithmic}
\end{algorithm}
